\documentclass[12pt,letterpaper, onecolumn]{exam}
\usepackage{graphicx}
\usepackage{caption}
\usepackage{color}
\usepackage{caption}
\usepackage{listings}
\usepackage[dvipsnames]{xcolor}
\graphicspath{ {./images/} }
\lhead{Mustafa Rashid\\}
\rhead{Homework 3: Relational DB Design\\}
\chead{\hline} 
\thispagestyle{empty} 
\newcommand*{\setdef}[1]{\left\{#1 \right\}} 
\renewcommand{\thepartno}{\Alph{partno}} % Force uppercase letters
\renewcommand{\partlabel}{\thepartno.}    % Format as A. B. C.

\begin{document}
\begingroup  
\noindent\LARGE CIS 5500: Database and Information Systems\\
\noindent\LARGE Homework 3: Transactions\\
\noindent\large \today\\
\noindent\large Mustafa Rashid\par
\noindent\large Spring 2026\par
\endgroup
\rule{\textwidth}{0.4pt}
\pointsdroppedatright
\printanswers
\renewcommand{\solutiontitle}{\noindent\textbf{Ans:}\enspace}

% SQL listing from https://gist.githubusercontent.com/vivngo/e37e1c7b79bf7e8b910c6566c59c46be/raw/db8723fc80c8b1b01665851106dba2f370b5d865/code.tex
\definecolor{dkgreen}{rgb}{0,0.6,0}
\definecolor{gray}{rgb}{0.5,0.5,0.5}
\definecolor{mauve}{rgb}{0.58,0,0.82}
\lstset{language=SQL,
  basicstyle={\small\ttfamily},
  belowskip=3mm,
  breakatwhitespace=true,
  breaklines=true,
  classoffset=0,
  columns=flexible,
  commentstyle=\color{dkgreen},
  framexleftmargin=0.25em,
  frameshape={}{}{}{}, %To remove to vertical lines on left, set `frameshape={}{}{}{}`
  keywordstyle=\color{blue},
  numbers=none, %If you want line numbers, set `numbers=left`
  numberstyle=\tiny\color{gray},
  showstringspaces=false,
  stringstyle=\color{mauve},
  tabsize=3,
  xleftmargin =1em
}

\begin{questions}
  \question \textbf{Question 1 (10 points)}
  \begin{solution}
    This is what is returned at timesteps 3-6. Timesteps 1 and 2 do not return anything.
    \begin{center}
\captionof{table}{Timestep 3}
\begin{tabular}{|c|c|c|c|}
\hline
\textbf{ID} & \textbf{Name} & \textbf{Zip} & \textbf{Balance} \\
\hline
210 & Madison & 19104 & 180 \\
\hline
\end{tabular}

\captionof{table}{Timestep 4}
\begin{tabular}{|c|c|c|c|}
\hline
\textbf{ID} & \textbf{Name} & \textbf{Zip} & \textbf{Balance} \\
\hline
410 & Luca & 19114 & 180 \\
\hline
\end{tabular}

\captionof{table}{Timestep 5}
\begin{tabular}{|c|c|c|c|}
\hline
\textbf{ID} & \textbf{Name} & \textbf{Zip} & \textbf{Balance} \\
\hline
210 & Madison & 19104 & 180 \\
310 & Noor & 19104 & 200 \\
\hline
\end{tabular}

\captionof{table}{Timestep 6}
\begin{tabular}{|c|c|c|c|}
\hline
\textbf{ID} & \textbf{Name} & \textbf{Zip} & \textbf{Balance} \\
\hline
210 & Madison & 19104 & 180 \\
410 & Luca & 19114 & 180 \\
310 & Noor & 19104 & 200 \\
\hline
\end{tabular}

\captionof{table}{Final State}
\begin{tabular}{|c|c|c|c|}
\hline
\textbf{ID} & \textbf{Name} & \textbf{Zip} & \textbf{Balance} \\
\hline
220 & Ethan & 19106 & 160\\
210 & Madison & 19104 & 180 \\
410 & Luca & 19114 & 180 \\
310 & Noor & 19104 & 200 \\
\hline
\end{tabular}
\end{center}
  \end{solution}
  \pagebreak
  \question \textbf{Question 2 (5 points)}
  \begin{solution}
     Postgres uses MVCC. Under \texttt{READ COMMITTED}, each statement gets a fresh snapshot of all committed data. This prevents dirty reads but allows phantom reads, since two statements in the same transaction can observe different committed states.
    At times 1 and 2 no rows qualify because the initial balances are below 180. At time 3, T1 sees its own uncommitted \texttt{UPDATE} to Madison. At time 4, T2 sees only its own uncommitted insert of Luca---T1's update is not yet committed, so it is invisible (no dirty read). At time 5, T1 sees both its own \texttt{UPDATE} and \texttt{INSERT}; T2's Luca is still uncommitted. At time 6, T1 has committed; T2's new statement snapshots the database and now sees T1's rows alongside its own Luca---a phantom read.
  \end{solution}
  \question \textbf{Question 3 (10 points)}
  \begin{solution}
    \begin{parts}
\part
      \textbf{What happens at time 5.}
      T1 attempts \texttt{INSERT INTO Accounts VALUES (410, 'Noor', 19104, 200)}. However, T2 already holds a write lock on \texttt{sid = 410} from its uncommitted insert of Luca (time 4). T1 therefore \emph{blocks} and waits for T2 to release its locks.

      \medskip
      \textbf{What happens at time 6 (T2 commits) and final state.}
      T2 executes its \texttt{SELECT} and then \texttt{COMMIT}s, releasing all of its locks. T1 resumes and attempts to complete the insert of \texttt{sid = 410}. Because T2 committed that row (Luca), a \emph{primary key conflict} arises. Postgres aborts T1's entire transaction with a unique constraint violation, rolling back both the \texttt{UPDATE} to Madison and the blocked insert of Noor. T2's \texttt{SELECT} at time 6 returns:

      \begin{center}
      \begin{tabular}{|c|c|c|c|}
      \hline
      \textbf{SID} & \textbf{Name} & \textbf{Zip} & \textbf{Balance} \\
      \hline
      410 & Luca & 19114 & 180 \\
      \hline
      \end{tabular}
      \end{center}

      The final state of Accounts is:

      \begin{center}
      \begin{tabular}{|c|c|c|c|}
      \hline
      \textbf{SID} & \textbf{Name} & \textbf{Zip} & \textbf{Balance} \\
      \hline
      220 & Ethan   & 19106 & 160 \\
      210 & Madison & 19104 & 90  \\
      410 & Luca    & 19114 & 180 \\
      \hline
      \end{tabular}
      \end{center}

      T2's committed data (Luca at \texttt{sid = 410}) persists. T1 was aborted, so Madison's balance reverts to 90 and Noor is never inserted.

      %% Part B
      \part
      \textbf{What happens at time 6 (T2 aborts) and final state.}
      T2's \texttt{ROLLBACK} undoes the insert of Luca (\texttt{sid = 410}) and releases all of T2's locks. T1, which was blocked waiting for that lock, can now proceed. T1's insert of \texttt{(410, 'Noor', 19104, 200)} succeeds because no row with \texttt{sid = 410} exists any longer. T1's \texttt{SELECT} returns:

      \begin{center}
      \begin{tabular}{|c|c|c|c|}
      \hline
      \textbf{SID} & \textbf{Name} & \textbf{Zip} & \textbf{Balance} \\
      \hline
      210 & Madison & 19104 & 180 \\
      410 & Noor    & 19104 & 200 \\
      \hline
      \end{tabular}
      \end{center}

      T1 then commits. The final state of Accounts is:

      \begin{center}
      \begin{tabular}{|c|c|c|c|}
      \hline
      \textbf{SID} & \textbf{Name} & \textbf{Zip} & \textbf{Balance} \\
      \hline
      220 & Ethan   & 19106 & 160 \\
      210 & Madison & 19104 & 180 \\
      410 & Noor    & 19104 & 200 \\
      \hline
      \end{tabular}
      \end{center}

      T2 rolled back, so Luca is never persisted. T1 successfully committed its \texttt{UPDATE} (Madison $\to$ 180) and its \texttt{INSERT} (Noor at \texttt{sid = 410}). Ethan is unchanged.
    \end{parts}
  \end{solution}
  \question \textbf{Question 4 (5 points)}
  \begin{solution}
    Under \texttt{REPEATABLE READ}, T2's snapshot is taken once at the start of its first statement (time 2). At that moment, no account satisfies \texttt{balance $\geq$ 180}. T1's \texttt{UPDATE} and \texttt{INSERT}s are committed at time 5---after T2's snapshot---so they remain invisible to T2's entire transaction. The only qualifying row T2 can see is the one it inserted itself: Luca. T2's \texttt{SELECT} at time 6 therefore returns:

    \begin{center}
    \begin{tabular}{|c|c|c|c|}
    \hline
    \textbf{SID} & \textbf{Name} & \textbf{Zip} & \textbf{Balance} \\
    \hline
    410 & Luca & 19114 & 180 \\
    \hline
    \end{tabular}
    \end{center}

    The final state of Accounts after both transactions commit is:

    \begin{center}
    \begin{tabular}{|c|c|c|c|}
    \hline
    \textbf{SID} & \textbf{Name} & \textbf{Zip} & \textbf{Balance} \\
    \hline
    220 & Ethan   & 19106 & 160 \\
    210 & Madison & 19104 & 180 \\
    410 & Luca    & 19114 & 180 \\
    310 & Noor    & 19104 & 200 \\
    \hline
    \end{tabular}
    \end{center}
  \end{solution}
  \pagebreak
  \question \textbf{Question 5 (5 points)}
  \begin{solution}
    Under \texttt{REPEATABLE READ}, Postgres takes a single MVCC snapshot at the transaction's first statement and reuses it for every subsequent statement. This prevents dirty reads, non-repeatable reads, and phantom reads.
    T2's snapshot was taken at time 2, when no account had balance $\geq$ 180. T1's writes were committed at time 5---after T2's snapshot---so they remain invisible to T2. Unlike \texttt{READ COMMITTED}, where T2's time-6 statement got a fresh snapshot and saw T1's commits, under \texttt{REPEATABLE READ} T2 only ever sees its own insert of Luca.
  \end{solution}
  \question \textbf{Question 6 (5 points)}
  \begin{solution}
    T2's \texttt{SELECT} at time 6 returns Madison, Luca, and Noor---identical to the \texttt{READ COMMITTED} result in Q1, and different from the \texttt{REPEATABLE READ} result in Q4 (which returned only Luca).

    Under \texttt{SERIALIZABLE}, READ/WRITE dependencies between transactions are tracked and a serialization anomaly leads to transaction failure (SSI). However, SSI only tracks dependencies between transactions \emph{both} running in \texttt{SERIALIZABLE}. Since T1 runs in \texttt{READ COMMITTED}, it is excluded from dependency tracking, so no serialization anomaly can be detected and T2 commits successfully. Because each query in T2 sees the committed version of the database as of the beginning of the transaction, and T1's writes were committed before T2's time-6 statement, T2 sees T1's rows. The effective serial order is T1 before T2.
  \end{solution}
  \question \textbf{Question 7 (10 points)}
  \begin{solution}
   At time 6, Postgres aborts T2 with \texttt{ERROR: could not serialize access due to read/write dependencies among transactions}.

    \begin{parts}
      \part
      In Q6, T1 ran in \texttt{READ COMMITTED}, so it was excluded from SSI's READ/WRITE dependency tracking. No serialization anomaly could be detected, and T2 committed successfully. The valid serial order was T1 before T2.

      In Q7, both transactions run in \texttt{SERIALIZABLE}, so READ/WRITE dependencies between them are tracked. The dependencies between T1 and T2 create a serialization anomaly with no valid serial order, leading to T2's failure.

      \part
      Under \texttt{SERIALIZABLE}, READ/WRITE dependencies between transactions are tracked and a serialization anomaly leads to transaction failure. Two rw-dependencies form a contradiction:

      \textbf{T2 $\xrightarrow{\text{rw}}$ T1:} T2 read \texttt{balance $\geq$ 180} at time 2 (empty). T1 later wrote rows satisfying that predicate (Madison $\to$ 180 at time 3, Noor at time 5). This requires T2 to be ordered before T1.

      \textbf{T1 $\xrightarrow{\text{rw}}$ T2:} T1 read \texttt{balance $\geq$ 180} at time 1 (empty). T2 later wrote a row satisfying it (Luca at time 4). This requires T1 to be ordered before T2.

      T1 must come before T2 \emph{and} T2 must come before T1 which is a contradiction. No valid serial order exists, so SSI aborts T2.
    \end{parts}
    \end{solution}
\end{questions}
\end{document}
